% MIN 0:00 à 0:20. Présenter le titre et le cadre général du stage.
\maketitle

% MIN 0:20 à 1:20. Situer le laboratoire et les rôles dans l'équipe.
\begin{frame}{Le CMAP et le groupe Robust Learning}
    \textbf{Centre de Mathématiques Appliquées (CMAP)}\\
    Laboratoire de l'École polytechnique.

    Le groupe \textbf{Robust Learning} étudie l'apprentissage robuste.

    \begin{itemize}
        \item \textbf{El-Mahdi El-Mhamdi} encadre les travaux scientifiques.
        \item \textbf{Peva Blanchard} fixe les priorités du produit.
        \item \textbf{Sébastien Rouault} assure la direction technique.
        \item \textbf{Mohamad-Ammar Said et moi} contribuons au développement.
    \end{itemize}

    Les contributions sont relues avant leur intégration, et les vérifications
    se font automatiquement.
\end{frame}

% MIN 1:20 à 2:00. Présenter le sujet général et le fil de l'exposé.
\begin{frame}{Le sujet du stage}
    J'ai participé au développement de \textbf{Krum}, une bibliothèque de
    composants logiciels pour mener des expériences en
    \textbf{apprentissage distribué robuste}.

    \begin{block}{Le problème étudié}
        Plusieurs machines entraînent un même modèle.
        Certains participants sont \alert{malhonnêtes} et faussent les calculs transmis.
    \end{block}

    \textbf{Plan}
    \begin{enumerate}
        \item Le problème et la mission.
        \item La construction de Krum.
        \item Les validations.
        \item Le bilan du stage.
    \end{enumerate}
\end{frame}
